% ==STYLE-BEGIN==
\documentclass[UTF8,11pt]{ctexart}
\usepackage{amsmath,amssymb,amsthm}
\usepackage[a4paper,margin=2.4cm,top=2.3cm,bottom=2.3cm]{geometry}
\usepackage{fancyhdr}
\usepackage{enumitem}
\usepackage{hyperref}
\hypersetup{hidelinks}

% ---- 版式 ----
\linespread{1.15}
\setlength{\parskip}{0.35em}
\setlength{\parindent}{1.6em}
\setlist{nosep,leftmargin=2em}

% ---- 页眉 ----
\pagestyle{fancy}
\fancyhf{}
\fancyhead[L]{\small\coursename}
\fancyhead[R]{\small\thepage}
\renewcommand{\headrulewidth}{0.4pt}

% ---- 定理环境：中文用正体，标题加粗，按小节共用一条流水号 ----
\newtheoremstyle{cnthm}{0.7em}{0.7em}{\normalfont}{}{\bfseries}{}{0.5em}{}
\theoremstyle{cnthm}
\newtheorem{thm}{定理}[section]
\newtheorem{lem}[thm]{引理}
\newtheorem{prop}[thm]{命题}
\newtheorem{cor}[thm]{推论}
\newtheorem{defn}[thm]{定义}
\newtheorem{rem}[thm]{注}
\newtheorem{exa}[thm]{例}
\newtheorem{conj}[thm]{猜想}
\renewcommand{\proofname}{\bfseries 证明}
\providecommand{\coursename}{}
% ==STYLE-END==
\renewcommand{\coursename}{解析几何\,·\,0928}
\title{几何与代数 0928}
\date{}
\begin{document}
\maketitle

\section{\ 向量的旋转}

$\vec e$ 是单位向量，向量 $\vec a$ 绕 $\vec e$ 旋转 $\theta$ 角，
得到 $\vec a'$。

\begin{enumerate}
\item 将 $\vec a$ 分解
\[
\vec a=\vec a_{\vec e}+\vec a_h,\qquad
\vec a_{\vec e}=(\vec a\cdot\vec e)\vec e,\qquad
\vec a_h=\vec a-\vec a_{\vec e}
=\vec a-(\vec a\cdot\vec e)\vec e .
\]
\item 将两向量 $\vec a_{\vec e},\vec a_h$ 分别绕 $\vec e$
旋转 $\theta$ 角，得到 $\vec a_{\vec e}',\vec a_h'$，有
\[
\vec a_{\vec e}'=\vec a_{\vec e},\qquad
\vec a_h'=\cos\theta\,\vec a_h+\sin\theta\,(\vec e\times\vec a_h).
\]
\item 合并
\[
\vec a_{\vec e}'=\vec a_{\vec e},\qquad
\vec a_h'=\cos\theta\cdot\vec a_h+\sin\theta\,(\vec e\times\vec a_h).
\]
\end{enumerate}

\subsubsection*{e.g.}

$\vec a=(1,1,1)$ 绕 $\vec e_1=(1,0,0)$ 右旋
$\dfrac\pi4$ 得到 $\vec a'$。

\begin{enumerate}
\item $\vec a=\vec a_{\vec e}+\vec a_h$，
$\vec a_{\vec e}=(1,0,0)$，$\vec a_h=(0,1,1)$；
\item
\[
\begin{aligned}
\vec a_h'
&=\cos\frac\pi4\vec a_h+\sin\frac\pi4(\vec e_1\times\vec a_h)\\
&=\left(0,\frac{\sqrt2}{2},\frac{\sqrt2}{2}\right)
+\frac{\sqrt2}{2}(0,-1,1)=(0,0,\sqrt2);
\end{aligned}
\]
\item $\vec a'=\vec a_h'+\vec a_{\vec e}=(1,0,\sqrt2)$。
\end{enumerate}

\section{\ 矩阵 (matrix)}

\textbf{定义}\quad $mn$ 个数排成 $m$ 行 $n$ 列的数表，
称为一个 $m\times n$（阶）矩阵。

例如 $(1,0,1)$ 是 $1\times3$ 矩阵；
$\begin{pmatrix}1\\0\\1\end{pmatrix}$ 是 $3\times1$ 矩阵。

矩阵一般用 $A,B,C,P,Q$ 表示。用 $[\ ]$ 或者 $(\ )$ 将数字括起来，
例如
\[
A=
\begin{pmatrix}
a_{11} & a_{12} & \cdots & a_{1n}\\
a_{21} & a_{22} & \cdots & a_{2n}\\
\vdots & \vdots & \ddots & \vdots\\
a_{m1} & a_{m2} & \cdots & a_{mn}
\end{pmatrix}
=(a_{ij})_{m\times n}.
\]
如果 $m=n$，则称该矩阵是一个 $n$ 阶方阵。如
$\begin{pmatrix}1&0&0\\0&1&0\\0&0&1\end{pmatrix}$
是一个 $3\times3$ 矩阵或三阶方阵。

\subsubsection*{矩阵的运算}

\begin{enumerate}
\item \textbf{加法}\quad $A=(a_{ij})$，$B=(b_{ij})$ 是
$m\times n$ 阶矩阵，则
\[
A+B\triangleq(a_{ij}+b_{ij})
\]
是 $m\times n$ 矩阵。
\item \textbf{数乘}\quad $k\in\mathbb R$，$A=(a_{ij})$ 是
$m\times n$ 矩阵，
\[
kA\triangleq(ka_{ij}).
\]
\item \textbf{乘法}

\textbf{(1) 矩阵与向量相乘}
\[
AX=
\begin{pmatrix}
a_{11} & a_{12} & \cdots & a_{1n}\\
a_{21} & a_{22} & \cdots & a_{2n}\\
\vdots & \vdots & \ddots & \vdots\\
a_{m1} & a_{m2} & \cdots & a_{mn}
\end{pmatrix}
\begin{pmatrix}
x_1\\x_2\\\vdots\\x_n
\end{pmatrix}
\triangleq
\begin{pmatrix}
a_{11}x_1+\cdots+a_{1n}x_n\\
a_{21}x_1+\cdots+a_{2n}x_n\\
\vdots\\
a_{m1}x_1+\cdots+a_{mn}x_n
\end{pmatrix}.
\]

\textbf{(2) 矩阵与矩阵相乘}\quad $A$ 是 $m\times n$ 矩阵，
$B$ 是 $n\times s$ 矩阵，
\[
\begin{aligned}
AB&=
\begin{pmatrix}
a_{11} & a_{12} & \cdots & a_{1n}\\
a_{21} & a_{22} & \cdots & a_{2n}\\
\vdots & \vdots & & \vdots\\
a_{m1} & a_{m2} & \cdots & a_{mn}
\end{pmatrix}
\begin{pmatrix}
b_{11} & b_{12} & \cdots & b_{1s}\\
b_{21} & b_{22} & \cdots & b_{2s}\\
\vdots & \vdots & & \vdots\\
b_{n1} & b_{n2} & \cdots & b_{ns}
\end{pmatrix}\\
&=
\begin{pmatrix}
\sum\limits_{k=1}^{n}a_{1k}b_{k1} & \cdots & \cdots
& \sum\limits_{k=1}^{n}a_{1k}b_{ks}\\
\vdots & & & \vdots\\
\vdots & & & \vdots\\
\sum\limits_{k=1}^{n}a_{mk}b_{k1} & \cdots & \cdots
& \sum\limits_{k=1}^{n}a_{mk}b_{ks}
\end{pmatrix}.
\end{aligned}
\]
\end{enumerate}

\end{document}
